\ifdefined\iscombined
	\lecturetitle{Environments}
\else
\documentclass{article}
\usepackage{listings}
\usepackage{amsthm}
\usepackage{comment}
\usepackage{amsmath}
\usepackage{amsfonts}
\usepackage{sectsty}
\usepackage{hyperref}
\usepackage{fancyhdr}
\usepackage[top=1in,bottom=1in,left=1in,right=1in]{geometry}
\usepackage{float}
\usepackage{graphicx}
\usepackage{tabularx}
\usepackage{mathtools}

\date{
	\vspace{-.75em}
	{\large \today}
}

\title{
	\vspace*{-1.5em}
	{\Huge Lecture 5} \\
	\vspace{-1em}
	\rule{\textwidth/4}{.01em} \\
	\vspace{-.25em}
	{\Large Environments}
	\vspace{-.75em}
}

\author{
	{\large Matthew Farah}
}

\hypersetup{
	colorlinks=true,
	linkcolor=blue,
	filecolor=magenta,
	urlcolor=blue,
	citecolor=blue,
	anchorcolor=blue
}

\fancyhead{}
\fancyhead[L]{\today}
\fancyhead[R]{Lecture 5 - Environments}

\sectionfont{\fontsize{12}{12}\bfseries}
\subsectionfont{\fontsize{10}{12}\normalfont}
\subsubsectionfont{\fontsize{10}{12}\normalfont}

\begin{document}

\maketitle
\pagestyle{fancy}

\fi

\begin{itemize}
	\item Systems are described by their design.
	\begin{itemize}
		\item Functional requirements provide the reason for the design's existence.
	\end{itemize}
	\item Systems always exist in some given environment.
	\begin{itemize}
		\item Non-functional requirements describe how the design should with its environment.
	\end{itemize}
	\item Systems and their environments are always interacting in some way.
	\begin{itemize}
		\item The environment always initiates these interactions.
	\end{itemize}
	\item A `scneario' is a sequence of interactions between a system and its environment.
	\begin{itemize}
		\item A global scenario is the integration of all scenarios together for some given system and environment pair.
	\end{itemize}
	\item A `business event' is something a system must perform when prompted by its environment.
	\begin{itemize}
		\item Scenarios terminate when a system fulfills the business task it was prompted to perform by its environment.
	\end{itemize}
	\item A `business process' is the act systems engage in to fulfil business events.
	\begin{itemize}
		\item One or more components of a system may be engaged in any given business process.
	\end{itemize}
	\item A `viewpoint' is a perspective that designs should take into account.
	\begin{itemize}
		\item For example, a legal viewpoint, an economic viewpoint, etc.
		\item Viewpoints are given for each business event.
		\begin{itemize}
			\item Inapplicable viewpoints should still be directly mentioned and specified as so.
		\end{itemize}
	\end{itemize}
	\item The modes of a system describe equivalence classes of similar system behaviours.
\end{itemize}

\ifdefined\iscombined
	\newpage
\else
	\end{document}
\fi
